\scanpage{80}
\mspage{86}

\deleted{b) Il faut prouver que si $p,p'$ sont des $\mathbb D$-fibrations, et $f$ une $\mathbb D$-équiv.\deleted{\illegible{1}}, alors $f$ induit des $\mathbb D$-équiv. sur les fibres. Considérons, pour $M\in\mathrm{Ob}\,A$, l'hom induit par $f$}
\deleted{\[\underset{\substack{\shortparallel\\H'}}{\underline{p'_* (M_{X'})}}\longrightarrow\underset{\substack{\shortparallel\\H}}{\underline{p_* (M_X)}}\,.\]}

\deleted{\deleted{\unsure{l'hypothèse|L'hypothèse}} Comme $f$ est une $\mathbb D$-fibration, c'est un hom d'objets loc. constants \inserted{sur $S$}, dont les fibres sont les $H^{\unsure{\bullet|\cdot|\prime}}_{\mathbb D}(X',M)$, $H^\bullet_{\mathbb D}(X,M)$ respectivement. Il faut prouver que l'hom. est iso, et tout ce qu'on sait, c'est qu'il induit $H^\bullet_{\mathbb D}(S,H')\unsure{-|\to}H^\bullet_S(S,H)$ qui est un iso. Si $H,H'$ étaient des \deleted{\unsure{fai|fais}} objets \uline{constants} de $\mathbb D(S)$, soient $P_S,Q_S$, alors}

Si $Y'=S$, l'énoncé \inserted{non prouvé} dit que si $f\colon X\longrightarrow S$ est une $\mathbb D$-fibration et une $\mathbb D$-équivalence, alors ses fibres sont $\mathbb D$-contractiles. Ce n'est pas sûr même pour une projection $X=T\times S\longrightarrow S$. Si $\unsure{y|\eta|\gamma|\mathfrak y}=\mathrm{hot}_0(\unsure{S|T}),\ \unsure{\zeta|z}=\mathrm{hot}_0(T)$, la question se réduit à celle-ci : Si $\forall M$ dans $A$, \deleted{\unsure{on a}}, l'hom
\deleted{\[(\zeta\times y)\times_{\mathbb D} M\xrightarrow{\sim}y\times_{\mathbb D} M\quad\text{alors}\]}
\[\alpha_M\unsure{=|:}\deleted{\unsure{y\times M|\zeta\times M}}\ \zeta\times_{\mathbb D}M\longrightarrow M\]
satisfait
\[y\times_{\mathbb D}\alpha_M\colon\ y\times_{\mathbb D}(\zeta\times_{\mathbb D} M)\xrightarrow{\sim}y\times_{\mathbb D} M\quad\text{\uline{iso}},\]
alors $\alpha_M$ lui-même est iso.
% archive sheet number: 79
